Congress possesses authority under Article I, \S4 to regulate the ``Manner'' of congressional elections, including the drawing of district lines. We analyze the Districting Integrity Act (DIA), a model federal statute requiring states to use recursive bisection to produce their congressional maps, and compare it against four alternative federal legislative designs: algorithm-specifying statutes, criteria-only statutes, commission-mandating statutes, and court-remedy statutes. The DIA's architecture is deliberately analogous to the Apportionment Act of 1941 (now codified at 2 U.S.C.\ \S2a), which mandated the Huntington-Hill method for apportioning seats among states. Like that statute, the DIA prescribes a specific deterministic algorithm---recursive bisection as implemented in the \texttt{bisect} platform---produces verifiable outputs, and preserves state responsibility for execution while preempting partisan manipulation. We assess each alternative design against four criteria: constitutional durability, political feasibility, VRA compatibility, and technical implementability. The DIA ranks highest overall. The criteria-only design (ranking second) provides a useful fallback if Congress lacks the technical expertise or political will to specify an algorithm. The commission-mandating and court-remedy designs are complements to, not substitutes for, an algorithm-based approach.
